%% kapa: exam=1 round=33 year=2022 period=1 form=A subject=civil_law subject_ko=민법 seq=5 number=5
%% source: Q-Net 감정평가사 기출문제 (cst00309) artlSeq=5209997 file=raw/1008/5209997/2022년도 제33회 감정평가사 1차시험 1교시(민법,경제학,부동산학).pdf
%% conversion: kordoc 4.24.51 -> tools/gen1.py (auto); review=proofread-sonnet; flags=none
\begin{question}{5}
민법상 법인에 관한 설명으로 옳지 않은 것은? (다툼이 있으면 판례에 따름)
\begin{choices}
  \item 재단법인의 정관변경은 그 정관에서 정한 방법에 따른 경우에도 주무관청의 허가를 얻지 않으면 효력이 없다.
  \item 사단법인과 어느 사원과의 관계사항을 의결하는 경우에는 원칙적으로 그 사원은 결의권이 없다.
  \item 사단법인의 사원자격의 득실에 관한 규정은 정관의 필요적 기재사항이다.
  \item 민법상 법인의 청산절차에 관한 규정에 반하는 합의에 의한 잔여재산처분행위는 특별한 사정이 없는 한 무효이다.
  \item 청산 중 법인의 청산인은 채권신고기간 내에는 채권자에 대하여 변제할 수 없으므로 법인은 그 기간 동안 지연배상 책임을 면한다.
\end{choices}
\end{question}
\answer{5}
